\section{Relation to earlier work}
\label{sec:history}

The entropy formula for Dicke tensors was proposed in the study of
asymptotic multipartite entanglement conversion. Vrana and Christandl
\cite[Section~VI, equation~(40)]{VC2015} suggested that the asymptotic
number of copies of a Dicke tensor needed per two-level GHZ tensor is
$1/\HH_2(p)$, where $p=\mathbf d/k$ and $\HH_2$ denotes entropy with
base-two logarithms. Equivalently, the logarithm of its asymptotic subrank
should be $\HH_2(p)$. Oneto and Ventura explicitly restated the ordinary
subrank formula as Conjecture~4.26 in \cite{OV2025}. Their symmetric-tensor
convention for a monomial differs from our Dicke tensor by a nonzero
scalar, which does not affect subrank. Permuting basis symbols also
identifies the partition convention used in some sources with arbitrary
positive multiplicity tuples. Throughout this paper, all $k$ tensor
factors remain separate; tensor powers group only corresponding factors.

A distinction between restriction and degeneration is needed to compare
our conclusion with Question~1.58 of Christandl, Vrana and Zuiddam
\cite{CVZ2019}. Our $\Qmon$ uses coordinate deletions followed by
independent relabellings and nonzero rescalings in the individual factors.
These are monomial restrictions in their Definition~1.6. Their monomial
subexponent $q_{\mathrm M}$, however, is defined using monomial
\emph{degeneration}: see Definition~1.7, equation~(1.9), and
Definition~1.18. Correspondingly, Proposition~1.19 expresses this
subexponent using monomial \emph{border} subrank. Remark~1.12 explicitly
distinguishes the resulting asymptotic rate from the rate defined using
monomial restrictions; the two are not identified there in general.
Writing superscripts to distinguish their notation, our theorem gives
\begin{equation}
\label{eq:cvz-convention-comparison}
\begin{aligned}
 \HH_2(p)
   &= \log_2 \Qmt(D_{\mathbf d}) \\
   &\le q_{\mathrm M}^{\mathrm{CVZ}}(D_{\mathbf d})
    \le q^{\mathrm{CVZ}}(D_{\mathbf d}) \\
   &= \log_2 \Qt(D_{\mathbf d})=\HH_2(p).
\end{aligned}
\end{equation}
The first comparison uses inclusion of monomial restriction in monomial
degeneration, and the second inclusion in ordinary degeneration. The
ordinary restriction--degeneration equivalence of
\cite[Theorem~1.11]{CVZ2019} justifies the last comparison with ordinary
asymptotic subrank. Thus the theorem also answers their Question~1.58,
without assuming any general equivalence between monomial restriction
and monomial degeneration. This external equivalence is needed only for
this comparison of conventions, not for our proof of the ordinary or
coordinate-restriction formula.

We use the higher-order Coppersmith--Winograd construction of
Christandl, Vrana and Zuiddam
\cite[Theorem~1.52 and Section~3.3]{CVZ2019}. The survival probability
of two colliding edges is
controlled by the rational rank of their coordinate differences. We use
this rank-dependent hashing mechanism and give the needed probability
calculation directly. The accompanying average-free sets also have a
standard antecedent: Vrana and Christandl prove the fixed-length
average-free estimate in \cite[Lemma~10 and Appendix~A]{VC2015}, quoted
as Proposition~3.2 in \cite{CVZ2019}. Our proof includes a no-carry digit
construction on a common Euclidean sphere. We make no novelty claim for
the hashing or average-free mechanisms.

Counting a specified class of candidates while testing isolation against
the entire surviving support likewise has a published precedent.
Vrana and Christandl's weighted isolation lemma
\cite[Lemma~2.3, arXiv version~2]{VC2019} permits an arbitrary nonnegative weight on the
full supported set and bounds its sum over globally isolated vertices
of the collision graph. An indicator weight selects a prescribed joint
type for counting purposes while leaving every possible collider in the
graph. Our argument uses this isolation principle with a fixed full
type and derives the required collision estimate from permutation
symmetry. The selection of a full joint type is not presumed to be a
local tensor restriction: the local restriction is obtained only after
isolation in the full support.

There is a related distinction in the entropy optimization of
\cite[Theorem~1.52]{CVZ2019}. Its feasible pair distributions prescribe
only the $2k$ one-coordinate marginals, not both full support-valued
marginal laws. An entropy bound proved for a uniform full first law
therefore cannot simply be substituted into that optimization. We do
not invoke the theorem as a black box or claim to evaluate its displayed
optimization. Instead, we count partners of each fixed candidate by
exact quotient fibers, and bound their entropy by an invariant-subspace
argument. This supplies the collision estimate needed for arbitrary
positive multiplicities, without assuming separate all-binary or
all-distinct cases.

Finally, diagonal slice rank and entropy upper bounds are established
parts of the slice-rank method; see Tao and Sawin
\cite[Example~5 and Proposition~6]{TaoSawin2016}. We include elementary
proofs of the facts used here. A score assigned to each local word
partitions every support entry into a slice and gives
$Q(D_{\mathbf d}^{\boxtimes N})\le k\exp(N\HH(p))$ for every $N\ge1$.
This upper bound applies to arbitrary complex linear maps and completes
the comparison with the coordinate restrictions constructed below.
